\section{Discussion}
\label{sec:discussion}

\paragraph*{What the results mean}
In this simulation, a behavioural signal that enters the offloading decision as a weighted term raises deadline success over the same rule without trust at every degraded fraction tested, and the node-local head comes within $0.001$--$0.003$ of an oracle that reads the sealed fault state (Sec.~\ref{sec:behav}).
% src: A5, A6
Most of that gain is already available from the node's own evidence: a non-learned monitor, H, which uses only the RSU's own \texttt{success\_ewma}, recovers most of the trust-aware gain, and the pre-registered rule did not separate it from the learned head (Sec.~\ref{sec:learning}).
% src: A27, A23 (X18: we do not say that learning adds nothing; the label is NULL under a 0.01 floor with signed differences 0.003-0.005)
We read this as a statement about the floor we fixed in advance and about this scenario, not as a statement that the learned head is unnecessary.
% src: X18

On graph structure, the load-bearing hypothesis was tested and is not supported.
The graph head did not beat the node-local head on the pre-registered diagnostic (NEGATIVE), in the closed loop, under staggered onset where a neighbour lead exists, or as a contained head that starts as C (NULL); and we measured why the first test could not have succeeded in the shared-onset scenario, as a boundary condition of that scenario (Sec.~\ref{sec:graph}).
% src: A9, A10, A11, A25, A28; X1, X2, X17, X21 (no generalisation: one injection model, S = 300, one scenario per seed, open-loop-trained heads, one D2 construction)
These results are for one injection model, one stagger size, one scenario per seed and heads trained on open-loop dispatch; we make no claim about other regimes or architectures.
% src: X17, X21

The weighted rule and a gate on the same head performed equivalently on task success (Sec.~\ref{sec:gate}).
The case for a weighted term is therefore a design argument: an exactly decomposable score and no empty-candidate-set failure mode, and not a measured gain.
% src: A15, X3
Layer 1 explains the score and is exact; it does not explain why a node's trust has its value.
Layer 2 attributes the trust head's output weakly, unstably and only asynchronously, and each Layer 2 explanation is one draw (Secs.~\ref{sec:layer1}, \ref{sec:layer2} and~\ref{sec:overhead}).
% src: A12, A13, A14; X15, X5

\paragraph*{Untested hypotheses}
The following were not run, and none of them is a finding.
They are the open questions that the record leaves, stated as hypotheses.
% src: A20 (verbatim list; "None of these, other than the single D2 construction of A28, was run and none may be reported as a finding")
\begin{itemize}
\item An observation-count feature.
\item Other architectures.
\item Training the trust heads on closed-loop dispatch, and the shift between the open-loop training dispatch and the closed-loop evaluation dispatch as a possible cause of D trailing C.
\item The self-supervised objective (next-task success of the node itself) rewards current behaviour only, so neighbour evidence that predicts a later failure may have no way to improve the objective (not tested).
\item Dilution of D's own-node signal by its aggregation over many neighbours: this was tested for one contained construction (D2, Sec.~\ref{sec:graph}), while dilution in general is untested.
\item Cold-start and collusion conditions where a graph model might help.
\item Gating a different head.
\end{itemize}
Staggered onset, which an earlier version of this list included, was tested (Sec.~\ref{sec:graph}).
% src: A20 (staggered onset tested in S13: A24, A25)
